% Phase 1 --- Object and Timing Audit
% Governed by: paper/spec/MAH_model_rebuild_execution_note_for_Codex.md (Sec. 4-6)
%              paper/spec/MAH_model_rebuild_amendment_v1.1_for_Codex.md (Sec. A, D, E)
% Status: Phase 1, APPROVED (replayed onto the correct base review-model@cd5b6e4;
% content unchanged from the original approval -- this module never referenced
% manuscript-specific objects, so no correction was needed here).
% Do not cite objects not yet defined in this module (RL-19); see
% paper/model_notes/01_symbols_and_objects.md for the deferred-objects list.

\section{Primitives, Objects, and Timing}
\label{sec:primitives-timing}

\subsection{Purpose}

This module fixes, before any payoff is derived, (i) a complete
classification of every primitive and state object used in the baseline,
and (ii) the timing of decisions. No payoff, cost, or value function is
introduced here. The advancement control $x_i$ is defined with the
narrower meaning required by amendment v1.1: it is \emph{not} generic R\&D,
upstream scientific research, or patent-generating effort.

\subsection{Object-status table}
\label{subsec:object-status}

Every symbol below is classified into exactly one of the nine categories
fixed by the execution note (primitive parameter; exogenous firm
characteristic; exogenous project characteristic; control; endogenous
firm-level object; endogenous route-level object; equilibrium price;
aggregate/distributional primitive; derived observed outcome).

\begin{center}
\begin{tabular}{lll}
\hline
Symbol & Category & Meaning \\
\hline
$M \in \{0,1\}$ & primitive parameter & institutional regime ($M=0$ pre-MAH, $M=1$ MAH) \\
$a_i>0$ & exogenous firm characteristic & research capability / research productivity \\
$k_i>0$ & exogenous firm characteristic & internal manufacturing capability \\
$z_j$ & exogenous firm characteristic & CMO manufacturing-service productivity \\
$q>0$ & exogenous project characteristic & project commercial/scientific value index \\
$m>0$ & exogenous project characteristic & manufacturing requirement / complexity \\
$g \in \{O,I\}$ & exogenous project characteristic & empirical novelty class (original / incremental) \\
$s(q)\in[0,1]$ & primitive (technology function of $q$) & downstream clinical/regulatory realization probability \\
$s_g(q)\in[0,1]$ & primitive (technology function of $q$) & type-specific realization probability, if required \\
$x_i \geq 0$ & control & project-development / advancement intensity \\
$H(a,k)$ & aggregate/distributional primitive & joint distribution of $(a_i,k_i)$ \\
$H_C(z)$ & aggregate/distributional primitive & distribution of CMO efficiency \\
$F(q,m)$ & aggregate/distributional primitive & distribution of project draws \\
$F_g(q,m)$ & aggregate/distributional primitive & type-specific distribution of project draws \\
$\lambda_i^{plan}=a_ix_i$ & endogenous firm-level object & arrival intensity of route-planning-stage projects \\
$r \in \{I,E,T,A\}$ & endogenous route-level object & commercialization route choice \\
$p_m$ & equilibrium price & manufacturing-service (CMO) price \\
\hline
\end{tabular}
\end{center}

No object outside this table is used in this module. Route values
$W_i^I,W_i^E,W_i^T,W^A$, the manufacturing-cost technologies, the expected
project value $\Omega_i$, and all equilibrium-market objects belong to
Phases 2--7 and are deliberately absent here (RL-19: no proposition or
payoff before its objects are formally defined).

\subsection{The advancement control $x_i$: narrowed interpretation}
\label{subsec:xi-narrow}

Per amendment v1.1 \S A, the baseline control $x_i$ must not be
interpreted as generic R\&D, upstream scientific research, or
patent-generating effort. It is instead
\[
x_i = \text{project-development / advancement intensity},
\]
or, more compactly, \emph{project-development intensity}. Its economic
role is to raise the flow of viable original-drug projects that reach the
route-planning / commercialization-relevant development stage:
\[
\lambda_i^{plan} = a_i x_i.
\]
Intended empirical counterparts: IND/application activity, initiation of
early-stage clinical development, first-stage clinical trials, and
advancement of viable drug candidates toward commercialization --- not raw
patent applications. Upstream scientific research, patent generation, and
discovery of basic compounds are \emph{outside} the baseline timing; they
may have occurred before a project enters this model. This model begins at
the margin where a developer decides how intensively to advance viable
drug projects while anticipating future commercialization routes, because
MAH directly changes downstream manufacturing/commercialization options,
not scientific-discovery technology.

Note on the pre-existing manuscript's terminology (see
\texttt{paper/manuscript/mah\_route\_indicator\_friction\_model.tex}
\S4.1, \S4.4): that manuscript's $x_i$ is described as ``original-drug
R\&D effort'' and its cost is written $C_i^R(x_i)=\frac{\kappa}{2}x_i^2$
with $R$ genuinely standing for R\&D. This is precisely the broad reading
amendment v1.1 requires narrowing (RL-21); the crosswalk from that
manuscript's $x_i,C_i^R$ to this module's $x_i,C_X$ is completed formally
in Phase 12, not here.

\subsection{Empirical novelty class $g$}
\label{subsec:novelty-class}

Per amendment v1.1 \S D, when the paper distinguishes original and
incremental innovation it uses an empirical project class $g\in\{O,I\}$.
This is \emph{not} a new baseline decision state: it enters only by
allowing project primitives to have type-specific distributions
$(q,m)\sim F_g(q,m)$ and, if required, a type-specific realization
probability $s_g(q)$. No object or proposition in this module ranks
$g=O$ against $g=I$; any such ranking is deferred to Phase 8's Corollary
and must be proved from primitive restrictions, not assumed.

\subsection{Timing}
\label{subsec:timing}

\paragraph{Stage 0: Institutional regime.} The economy is under
$M\in\{0,1\}$. $M=0$ is the pre-MAH organizational environment; $M=1$ is
the regime in which MAH retained holder--producer separation is legally
available. (The pre-existing manuscript additionally uses a continuous
implementation index $\eta\in[0,1]$ conditional on $M=1$; per RL-13, this
module and the new baseline do not use $\eta$ at all.)

\paragraph{Stage 1: Project-development / advancement investment.}
(Relabeled per amendment v1.1 \S E; was ``R\&D'' in the pre-amendment
draft.) Developer $i$ observes $(a_i,k_i)$ and chooses
\[
x_i \geq 0.
\]
The intensity of viable projects reaching the route-planning /
commercialization-relevant development stage is
\[
\lambda_i^{plan} = a_i x_i.
\]
(This production function is mathematically identical to the pre-existing
manuscript's equation~(eq:planning\_arrival); only the interpretation of
$x_i$ narrows.)

\paragraph{Stage 2: Project characteristics.} Each arriving project draws
\[
(q,m) \sim F(q,m),
\]
optionally stratified by novelty class, $(q,m)\sim F_g(q,m)$. (The
pre-existing manuscript has no analog of this stage: its route payoffs
$G_i^I,G_i^E,G_i^T,G_i^A$ are firm-level constants, not functions of a
project draw $(q,m)$. This is new structure the rebuild adds.)

\paragraph{Stage 3: Commercialization organization.} Conditional on the
project draw and the equilibrium manufacturing-service price $p_m$, the
project chooses among
\[
\{I,E,T,A\},
\]
where $I$ is internal manufacturing, $E$ is retained entrusted
manufacturing, $T$ is transfer/out-license/non-retained route, and $A$ is
abandonment or indefinite delay. Route \emph{values} are not defined in
this module; only the choice set and its timing are fixed here. (The
pre-existing manuscript resolves this choice via a logit model over
$G_i^r$, eq.~(eq:logit\_prob); this module and Phase 4 use a deterministic
maximum instead, per RL-12 -- logit is deferred to Phase 11 as an
explicit, non-baseline calibration extension.)

\paragraph{Stage 4: Downstream realization.} A route-independent
downstream clinical/regulatory realization probability $s(q)\in[0,1]$ (or
$s_g(q)$ if type-specific) determines whether the retained project
generates commercial value. $M$ does not directly shift $s(q)$ or
$s_g(q)$ (RL-09). (The pre-existing manuscript's analog is
$\zeta_i^r=s_i\chi_i^r$, a firm-level rather than project-level object;
crosswalked formally in Phase 12.)

\paragraph{Stage 5: Manufacturing-service market consistency.} Aggregate
entrusted-route demand must be consistent with the equilibrium CMO price
$p_m^*$. Market clearing itself is defined in Phase 6--7. (The
pre-existing manuscript's analog is $p_m^*(M,\eta)$, Proposition
prop:cmo\_existence; the new baseline drops the $\eta$ argument per
RL-13.)

\subsection{What this module deliberately excludes}

No demand curve, no profit function, no manufacturing-cost technology, no
route value, no advancement cost, no expected project value $\Omega_i$,
and no CMO supply/demand object appears in this module. No claim is made
here about whether MAH raises patent applications, upstream research, or
breakthrough innovation (RL-21--RL-24): this module only fixes objects and
timing. No logit, inclusive value, or continuous $\eta$ appears anywhere
in this module (RL-12, RL-13).
